\chapter{Empirical Results \& Analysis}
\label{ch:results_simple}

This chapter reports every experiment on the simple-modelling track in the order the methodology builds: Phase~0 validates the learning algorithms on games with known outcomes; Phase~1 applies the validated PPO loop to real SPY data under Iteration~1 (3-D state) and Iteration~2 (5-D state). All numbers are produced by committed scripts; none are illustrative placeholders.

\section{Experimental setup and evaluation metrics}
\label{sec:setup_simple}

\paragraph{Phase 0 (games).} Fixed seed 42. Budgets: 50{,}000 environment timesteps on CartPole; 150{,}000 on Flappy Bird and LunarLander; REINFORCE uses an equivalent episode budget. Evaluation: 30 greedy episodes. Metric: mean evaluation return.

\paragraph{Phase 1 (trading).} SPY daily closes 2018--2024, chronological 58 / 26 month split (test $=$ Nov 2022 -- Dec 2024). $C_0=\$10{,}000$, $c=5$\,bps, PPO 80{,}000 timesteps, seed 42. Metrics: wealth change $\Delta W$ (\$), trades per month, win rate vs the strongest passive baseline. Standing diagnostics: accounting identity (tolerance $10^{-6}$) and illegal-action request rate.

\section{Phase 0: algorithm validation on games}
\label{sec:phase0_results}

Table~\ref{tab:phase0_results} and Figures~\ref{fig:phase0_cartpole_bar}--\ref{fig:phase0_lunar_bar} summarise the three-game suite.

\begin{table}[h]
\centering
\caption{Phase 0 mean evaluation return (30 greedy episodes, seed 42). CartPole solved $\approx 475$; LunarLander $\approx 200$. Tabular Q-learning only on discretized CartPole.}
\label{tab:phase0_results}
\begin{tabular}{lrrr}
\toprule
\textbf{Algorithm} & \textbf{CartPole-v1} & \textbf{FlappyBird-v0} & \textbf{LunarLander-v3} \\
\midrule
Random                  & 28.8           & $-7.4$        & $-183.4$ \\
Tabular Q-learning      & \textbf{500.0} & ---           & --- \\
DQN                     & \textbf{500.0} & 7.0           & \textbf{177.6} \\
REINFORCE (scratch)     & 292.2          & 7.1           & 15.3 \\
A2C                     & \textbf{500.0} & 4.6           & $-41.2$ \\
PPO                     & \textbf{500.0} & \textbf{12.6} & 176.4 \\
\bottomrule
\end{tabular}
\end{table}

\begin{figure}[h]
  \centering
  \includegraphics[width=0.85\textwidth]{phase0_cartpole_comparison}
  \caption{CartPole algorithm comparison. Learning methods clear the random floor; Q-table / DQN / A2C / PPO hit the 500-step cap.}
  \label{fig:phase0_cartpole_bar}
\end{figure}

\begin{figure}[h]
  \centering
  \includegraphics[width=0.85\textwidth]{phase0_cartpole_curves}
  \caption{CartPole learning curves (25-episode moving average). Scratch REINFORCE learns but plateaus below the clipped / actor--critic methods.}
  \label{fig:phase0_curves}
\end{figure}

\begin{figure}[h]
  \centering
  \includegraphics[width=0.85\textwidth]{phase0_flappy_comparison}
  \caption{Flappy Bird comparison. PPO is the strongest learner under the fixed budget.}
  \label{fig:phase0_flappy_bar}
\end{figure}

\begin{figure}[h]
  \centering
  \includegraphics[width=0.85\textwidth]{phase0_lunarlander_comparison}
  \caption{LunarLander-v3 comparison. PPO and DQN approach the solved threshold ($\approx 200$).}
  \label{fig:phase0_lunar_bar}
\end{figure}

\paragraph{Carry-forward.} Every implementation clears its random floor. Scratch REINFORCE works but is noisy --- the textbook motivation for actor--critic and clipping. Tabular Q-learning reaches 500 on CartPole \emph{only after} hand-binning; trading states with continuous cash and prices make that table explode, so Phase~1 uses policy-gradient methods with function approximation. PPO is the most consistent learner across games and is therefore the trading algorithm.

\section{Phase 1, Iteration 1: the 3-D state agent}
\label{sec:phase1_iter1_results}

Table~\ref{tab:phase1_iter1} and Figure~\ref{fig:phase1_iter1_bar} compare PPO on \texttt{TradingEnv-v0} against four baselines with explicit objectives.

\begin{table}[h]
\centering
\caption{Iteration 1 on 26 held-out SPY test months. $\Delta W$ in dollars; fee 5\,bps; PPO 80k steps.}
\label{tab:phase1_iter1}
\begin{tabular}{lrrrr}
\toprule
\textbf{Method} & \textbf{Mean $\Delta W$} & \textbf{Std} & \textbf{Trades/mo} & \textbf{Win vs B1b} \\
\midrule
B0 do-nothing        & 0.00   & 0.00   & 0.0  & 31\% \\
B1 buy-one-hold      & 7.97   & 16.67  & 1.0  & 31\% \\
B1b buy-max-hold     & 76.91  & 178.68 & 19.1 & --- \\
B2 random (masked)   & 14.14  & 32.93  & 12.8 & 31\% \\
A1 PPO (3-D state)   & 76.91  & 178.68 & 19.1 & 0\% \\
\bottomrule
\end{tabular}
\end{table}

\begin{figure}[h]
  \centering
  \includegraphics[width=0.85\textwidth]{phase1_spy_daily_delta_w}
  \caption{Iteration 1 mean $\Delta W$ on test months. A1 matches B1b exactly once the fair fully-invested baseline is included.}
  \label{fig:phase1_iter1_bar}
\end{figure}

\begin{figure}[h]
  \centering
  \includegraphics[width=0.85\textwidth]{phase1_spy_daily_wealth_path}
  \caption{Wealth path for a representative test month: B1b buy-max hold vs A1 PPO (identical after collapse).}
  \label{fig:phase1_wealth}
\end{figure}

\paragraph{Honest headline.} Deterministic PPO matched B1b on \textbf{all 26} test months: under pure $\Delta W$ with a 3-D state, the learned policy is ``get fully invested and stay there.'' Against the one-share baseline B1 this would look like a $10\times$ win; against fair B1b it is rediscovered buy-and-hold. Accounting identity held exactly; illegal-action request rate was 3.9\%.

\section{Phase 1, Iteration 2: the 5-D state agent}
\label{sec:phase1_iter2_results}

Iteration 2 re-trains the identical PPO configuration on \texttt{TradingEnv-v1}, changing only the state. Table~\ref{tab:phase1_iter2} and Figure~\ref{fig:phase1_iter2_bar} report the comparison.

\begin{table}[h]
\centering
\caption{Iteration 2 vs Iteration 1 on the same 26 test months, budget, fee and seed.}
\label{tab:phase1_iter2}
\begin{tabular}{lrrr}
\toprule
\textbf{Method} & \textbf{Mean $\Delta W$} & \textbf{Trades/mo} & \textbf{Months identical to B1b} \\
\midrule
A1 PPO, 3-D (Iter.\,1) & 76.91 & 19.1 & 26 / 26 \\
A2 PPO, 5-D (Iter.\,2) & 0.00  & 0.0  & 0 / 26 \\
B1b buy-max-hold       & 76.91 & 19.1 & --- \\
\bottomrule
\end{tabular}
\end{table}

\begin{figure}[h]
  \centering
  \includegraphics[width=0.85\textwidth]{phase1_iteration2_delta_w}
  \caption{Iteration 2 (5-D) vs Iteration 1 / buy-max. Richer state breaks the collapse but converges mostly-flat under deterministic evaluation.}
  \label{fig:phase1_iter2_bar}
\end{figure}

The richer state \emph{did} break the buy-and-hold collapse (0/26 months identical to B1b) but converged to the opposite corner: a cautious, mostly-flat policy. Mean action probabilities on a representative test month: Hold 0.66, Buy 0.12, Sell 0.22 --- so argmax evaluation never trades. Sampling actions stochastically yields mean $\Delta W = +\$1.87$. Month-by-month, the flat agent still beats fully-invested B1b in 8 of 26 months (the falling ones).

\section{Master comparison across iterations}
\label{sec:master_table}

\begin{table}[h]
\centering
\caption{Master results table for the simple-modelling track (Phase 1, SPY test months).}
\label{tab:master}
\begin{tabular}{lccc}
\toprule
\textbf{Metric} & \textbf{B1b buy-max} & \textbf{Iter.\,1 PPO (3-D)} & \textbf{Iter.\,2 PPO (5-D)} \\
\midrule
State dimension              & ---     & 3       & 5 \\
Mean $\Delta W$ (\$)         & 76.91   & 76.91   & 0.00 \\
Mean trades / month          & 19.1    & 19.1    & 0.0 \\
Identical to B1b (months)    & ---     & 26/26   & 0/26 \\
Accounting identity          & pass    & pass    & pass \\
\bottomrule
\end{tabular}
\end{table}

\section{Analysis and key findings}
\label{sec:analysis_simple}

\paragraph{1. The reward, not the state, sets the incentive.} Under $r_t=\Delta W_t$ in a mostly-rising test window, full investment is near-optimal and Iteration~1 found exactly that. Adding entry and time context moved the policy along the exposure axis (always-invested $\to$ almost-never-invested) without creating selective timing. Both extremes are rational responses to a risk-neutral reward estimated from noisy monthly episodes whose training window includes the COVID crash and the 2022 bear market.

\paragraph{2. Baseline design is itself a result.} Without fair B1b, Iteration~1 would have been reported as ``PPO beats buy-and-hold by $10\times$.'' With it, the correct statement is ``PPO rediscovered buy-and-hold.''

\paragraph{3. What Iteration 3 must change.} Because the state upgrade only shifted exposure, the next iteration targets the objective: $r_t = \Delta W_t - \lambda D_t$ (drawdown penalty) gives a reason to hold \emph{selectively}. The masked-action rate and accounting identity carry forward as standing diagnostics.

\paragraph{Scope.} Single-seed, fixed budget, monthly daily-bar episodes. Multi-seed confidence intervals, minute-bar episodes ($T{=}390$), and the drawdown-penalised reward are the planned extensions, in that order.
